\section{Exit the Cave}

Plato\footnote{\url{https://www.hermes-press.com/life_awakening.htm}}:

\begin{quotationx}
In the \emph{Theaetetus}, Socrates asks: ``How can you determine whether at this moment we are sleeping, and all our thoughts are a dream; or whether we are awake, and talking to one another in the waking state?''

Theaetetus: ``Indeed, Socrates, I do not know how to prove the one any more than the other, for in both cases the facts precisely correspond; and there is no difficulty in supposing that during all this discussion we have been talking to one another in a dream; and when in a dream we seem to be narrating dreams, the resemblance of the two states is quite astonishing.''

Socrates: ``You see, then, that a doubt about the reality of sense is easily raised, since there may even be a doubt whether we are awake or in a dream. And as our time is equally divided between sleeping and waking, in either sphere of existence the soul contends that the thoughts which are present to our minds at the time are true; and during one half of our lives we affirm the truth of the one, and, during the other half, of the other; and are equally confident of both.''

\end{quotationx}


The Church:

\begin{quotex}
For this reason it says, ``Awake, sleeper, And arise from the dead, And Christ will shine on you.'' \flright{\textsc{Ephesians 5:14}}
\end{quotex}


Evola\footnote{\url{http://www.gornahoor.net/PaganImperialism/ThoseWhoKnow.htm}}:

\begin{quotex}
The sacred and the divine are matters of faith. This is the truth which has been imposed on Europe of late. Our truth is otherwise: it is better to know that we don't know rather than to believe. In the contemporary mentality, there is a central point at which the attitudes of materialistic science and religion meet: in an identical renunciation, in an identical pessimism, in an identical agnosticism about the spiritual, declared and methodical in one case, veiled in the other. The premise of materialistic science is basically that science---in the sense of real, positive and empirical knowledge---can only subsist in what is physical; and that in the non-physical there can be no science, so that the scientific method neglects it and abandons it, by lack of authority, to belief, to the dull and arbitrary abstractions of philosophy, or to the ``exigencies'' of sentiment and morality. In addition, religion, insofar as it is focused exclusively on faith and does not admit an esoteric initiatory teaching beyond the profane religion imposed on the masses, or a gnosis beyond pious superstition, ends up with the same renunciation. In fact, one believes only when one does not know and thinks one cannot know. Hence, there is again the same agnosticism of the ``positivists'' with respect to whatever is not material and gross reality.
\end{quotex}

The first step is a sort of ``waking up'' -- Evola was writing after centuries of accretion and decay, in the dawn of the modern era, and he saw how faith (pistis) separated from knowledge (gnosis) would simply be another form of numbness or deadness to what is real.

Traditionalists of every stripe, as well as conservatives, wish to ``stem the tide'', either politically or culturally or theologically. Evola draws attention to their errors in an essay\footnote{\url{http://www.new-antaios.net/2011/04/julius-evola-on-ernst-junger-east-and-west-\%E2\%80\%93-the-gordian-knot/}} on Junger's work:

\begin{quotex}
The merit of Jünger in that first phase of his thought is that he had rec\-og\-nized the fatal error of those who think that every\-thing may be brought back to order, that this new men\-ac\-ing world, ever advanc\-ing, may be sub\-dued or held on the basis of the vision of life of the val\-ues of the pro\-ceed\-ing age, that is to say of bour\-geois civ\-i\-liza\-tion. If a spir\-i\-tual cat\-a\-stro\-phe is to be averted mod\-ern man must make him\-self capa\-ble of devel\-op\-ing his own being in a higher dimen\-sion -- and it is in this con\-nex\-ion that Jünger had announced the above-mentioned watch\-word of ``heroic real\-ism'' and pointed out the ideal of the ``absolute per\-son,'' capa\-ble of mea\-sur\-ing him\-self with ele\-men\-tary forces, capa\-ble of seiz\-ing the high\-est mean\-ing of exis\-tence in the most destruc\-tive expe\-ri\-ences, in those actions wherein the human indi\-vid\-ual no longer counts... Now it is impor\-tant to point out that wher\-ever forces belong\-ing to the first of the three domains emerge and break forth, only the pos\-si\-bil\-i\-ties of the third domain can really face them. Any attempt to stem on the basis of forces and val\-ues of the inter\-me\-di\-ate zone can only be pre\-car\-i\-ous, pro\-vi\-sional and rel\-a\-tive.
\end{quotex}


Evola's observations could be equally taken to apply in the cultural or religious sphere, as well as political, perhaps even more so. The ``answer'' to the ``Decline of the West'' can only come in the form of regeneration for the person or individual, outside of all structure if need be, church or state. In contrast to the seamless teaching of ancient wisdom (which continued in some consciousness down through our own day, as Evola attempted to enumerate), the modern solution is very much in the spirit of Kant, who ``destroyed knowledge in order to make room for faith''. Rather than oppose the Zeitgeist, it divinizes it.

One doesn't have to have the insight of B. Mouravieff to see that the Zeitgeist is in some sense ``winning''. As Rosenstock-Huessy notes, the modern world has melted everything but the atomistic individual and the Leviathan state. We cannot regenerate the Beast. At least, we do not start there. It is the individual who must ``skillfully use'' (and be ``skillfully used''). The era of the individual is here:

\begin{quotex}
Today we are living through the agonies of transition to the third epoch. We have yet to establish Man, the great singular of humanity, in one household, over the plurality of races, classes, and and age groups. This will be the center of struggle in the future. They pose the questions the Third Millennium will have to answer... the State is on the defensive because it is inadequate for the needs of the coming age. The theme of future history will be not territorial nor political, but social...

Huessy's\footnote{\url{http://en.wikipedia.org/wiki/Eugen\_Rosenstock-Huessy}} sequence ran:

\begin{enumerate}


\item The Christianization of Europe/Romanization of the Church

\item The Papal Revolution 11\textsuperscript{th}-13\textsuperscript{th} century, where the Church destroys the yoke of State

\item Protestant Reformation

\item The Enlightenment
\end{enumerate}
\end{quotex}


Huessy's flawed prophecy lacked any answer, but he saw what was upon us-

``Say neither it is blessed, nor it is cursed, but only `it is here'... ''

What we see now is a desperate confusion, very far indeed from the harmony of true sight into Reality, which was creative, synthesizing, and eclectic, and managed to reconcile both At hens \& Jerusalem.

As Gornahoor has noted, reinventing \emph{ad hoc} answers or worldviews (in the modern jargon, ``critical studies'') is essentially more of the problem, not a solution. It is not ``riding the tiger''. It is watching the shadows of the puppets on the wall, and convincing others to do the same, while pretending that one is not in a cave, and that this is all that there is.

The ``excluded middle'' of partial ignorance does not exist in this atmosphere either. As conservative inability to mount opposition (either externally or internally) becomes more clear, there will become more people convinced that radical regeneration along ancient lines is the only possible exit from the Cave.


\flrightit{Posted on 2011-05-20 by Logres}

